\documentclass[11pt,a4paper]{article}
\usepackage[margin=27mm]{geometry}
\usepackage{amsmath,amssymb,amsthm,hyperref}
\newtheorem{proposition}{Proposition}
\newcommand{\gen}[1]{\langle#1\rangle}
\newcommand{\Sz}{\operatorname{Sz}}
\title{Panels of eight retained Suzuki markings}
\author{Human Compression project}\date{}
\begin{document}\maketitle
\begin{abstract}
We verify the actual incidence panels of eight previously constructed
markings of $\Sz(8)$. Together with the previously reviewed ninth marking,
these give thin, flag-transitive, residually connected geometries for
the nine indicated parameter pairs. This finite result does not assert
that the markings exhaust the possible embeddings or prove the general
thinness assertion in the source.
\end{abstract}
\section{Matrices and the finite certificate}
Use $\mathbb F_8=\mathbb F_2[X]/(X^3+X+1)$, with integers $0,\ldots,7$
encoding the coefficients in the basis $1,X,X^2$. Thus the integer $2$
denotes $X$, and addition is exclusive-or. In the classical Suzuki model put
\[
 w=\begin{pmatrix}0&0&0&1\\0&0&1&0\\0&1&0&0\\1&0&0&0\end{pmatrix},\qquad
 z(t)=\begin{pmatrix}1&0&0&0\\0&1&0&0\\t&0&1&0\\t^4&t&0&1\end{pmatrix},
 \quad h(t)=\operatorname{diag}(t^5,t,t^{-1},t^{-5}).
\]
We use the classical ambient matrix model and the classification of its
maximal subgroups as explicit premises. The certificate checks the matrices,
subgroups and cosets below directly, without enumerating the ambient group.
The retained recipes, in their original order, are
\[
\begin{array}{c|rrrrrrrr}
(p,p')&(2,5)&(2,7)&(2,13)&(5,7)&(5,13)&(7,7)&(7,13)&(13,13)\\
x&1&4&3&1&1&4&2&3\\
y&-&-&-&2&3&-&3&-
\end{array}
\]
For $p=2$, take $(a,b,c)=(w,z(2),z(x))$ and
$(G_0,G_1,G_2)=(\gen{a,b},\gen{b,c},\gen{a,c})$.
For $p=p'$ odd, take $(a,b,c)=(w,z(2)wz(2),z(x))$.
In the remaining cases put $\lambda=(y^4x^{-4})^3$ and take
$(a,b,c)=(w,h(\lambda)w,z(x))$.
For both odd cases use
$(G_0,G_1,G_2)=(\gen{a,b},\gen{a,c},\gen{b,c})$.

The returned lists, checked independently by ordinary finite-field
multiplication, have orders $14,2p,2p'$; all three pair intersections
have order two and their triple intersection is trivial. The reader checks
every product and inverse in each listed subgroup, and verifies that its
two specified involutions generate the whole list. These data in
particular identify the subgroups as the stated dihedral groups, with
$D_4=C_2\times C_2$ when $p=2$.

\section{From the returned cosets to the geometry}
For missing type $i$ and other types $j,k$, list all distinct left cosets
$gG_i$ with $g\in G_j$, retaining precisely those that meet $G_k$.
Every type-$i$ object incident with $G_j$ has a representative in $G_j$,
so this procedure exhausts the actual panel above the edge $(G_j,G_k)$.
It uses at most thirteen candidates for any one panel. The certificate
retains the entire coset and its intersections with $G_j$, $G_k$, and
$G_j\cap G_k$, not only a claimed cardinality.

For each of the eight recipes and each of the three missing types, the
independent reader obtains exactly two cosets, both meeting
$G_j\cap G_k$. There are thus twenty-four panels of size two, and every
chamber containing a base edge has a common representative.

\begin{proposition}
Each of the eight coset geometries is thin, flag transitive and residually
connected. It satisfies RWPRI and local $2$-transitivity under the literal
algebraic-residue convention of the fixed source.
\end{proposition}
\begin{proof}
Every incident pair has a common representative, so it translates to a
base edge. The calculation therefore proves that every panel has size
two. It also proves that every chamber translates to the base chamber:
after translating one of its edges to $(G_j,G_k)$, the remaining coset
meets $G_j\cap G_k$. This proves flag transitivity; every flag extends
to a chamber by the common-representative argument in rank three.

Each $G_i$ is generated by its two distinct order-two intersections
with the other parabolics. Thus its rank-two residue is connected:
successive multiplication by the two intersection subgroups gives an
incidence walk to every coset. The whole incidence graph is connected
as well. Indeed $G_0=D_{14}$ is maximal in $\Sz(8)$, and an additional
parabolic has intersection of order two with $G_0$, so the three
parabolics generate the ambient group. The coset-walk argument applies
again. Rank-one residues impose no further connectedness requirement.

For the empty flag, $G_0$ supplies the required maximal subgroup. In
each parabolic, either of its order-two intersection subgroups is maximal,
since its index is one of $2,5,7,13$. For a base edge the trivial triple
intersection is maximal in its order-two stabilizer. These are all the
RWPRI conditions. Finally the algebraic rank-one residues are the regular
two-point actions of the pair intersections, and hence are $2$-transitive.
\end{proof}

\paragraph{Scope and attribution.}
The constructions and the earlier $(5,5)$ example are retained phase-five
work; the latter already has a distinct accepted panel certificate.
The new contribution here is the complete literal incidence check and
its application for the other eight saved markings. The statement concerns
these markings, not all embeddings with the same orders. The source is
D.~Leemans, \emph{The rank 3 geometries of the simple Suzuki groups
$\Sz(q)$}, Note di Matematica 19(1) (1999), 43--63, especially its
definitions on printed pp.~45--46 and Theorem~6.1 on p.~61.
The source's algebraic residue definition and actual thinness conclusion
remain distinct notions until proved to agree in an application, as above.
There is no new source-completion, refutation, conjecture, Lean or
compression credit from these positive checks.
\end{document}
